\documentclass[11pt,fleqn]{article}
\usepackage[margin=1in,top=1in,bottom=1in]{geometry}
\usepackage{tikz}
\usepackage{mathtools}
\usepackage{longtable}
\usepackage{enumitem}
\usepackage[hidelinks]{hyperref}
\usepackage[natbibapa]{apacite}
%\usepackage[dvips]{graphics}
%\usepackage[table]{xcolor}
%\usepackage{amssymb}
\usepackage{float}
%\usepackage{subfig}
\usepackage{booktabs}
\usepackage{subcaption}
\usepackage{colortbl}
\usepackage{color}

\usepackage[normalem]{ulem}

\usepackage{multicol}
\usepackage{txfonts}
\usepackage{amsfonts}
%\usepackage{natbib}
\usepackage{gb4e}
\usepackage[all]{xy}
\usepackage{rotating}
\usepackage{tipa}
\usepackage{multirow}
\usepackage{authblk}
\usepackage{url}
\usepackage{pdflscape}
\usepackage{rotating}
\usepackage{adjustbox}
\usepackage{array}


\def\bad{{\leavevmode\llap{*}}}
\def\marginal{{\leavevmode\llap{?}}}
\def\verymarginal{{\leavevmode\llap{??}}}
\def\swmarginal{{\leavevmode\llap{4}}}
\def\infelic{{\leavevmode\llap{\#}}}

\definecolor{airforceblue}{rgb}{0.36, 0.54, 0.66}
%\definecolor{gray}{rgb}{0.36, 0.54, 0.66}

\definecolor{Pink}{RGB}{240,0,120}
\newcommand{\red}[1]{\textcolor{Pink}{#1}}
\newcommand{\jd}[1]{\textbf{\textcolor{Pink}{[jd: #1]}}}
\definecolor{green}{RGB}{0,158,115}
\definecolor{orange}{RGB}{213,94,0}

\newcommand{\dashrule}[1][black]{%
  \color{#1}\rule[\dimexpr.5ex-.2pt]{4pt}{.4pt}\xleaders\hbox{\rule{4pt}{0pt}\rule[\dimexpr.5ex-.2pt]{4pt}{.4pt}}\hfill\kern0pt%
}

\setlength{\parindent}{.3in}
\setlength{\parskip}{0ex}

\newcommand{\yi}{\'{\symbol{16}}}
\newcommand{\nasi}{\~{\symbol{16}}}
\newcommand{\hina}{h\nasi na}
\newcommand{\ina}{\nasi na}

\newcommand{\foc}{$_{\mbox{\small F}}$}

\hyphenation{par-ti-ci-pa-tion}

\setlength{\bibhang}{0.5in}
\setlength{\bibsep}{0mm}
\bibpunct[:]{(}{)}{;}{a}{}{,}

\newcommand{\6}{\mbox{$[\hspace*{-.6mm}[$}} 
\newcommand{\9}{\mbox{$]\hspace*{-.6mm}]$}}
\newcommand{\sem}[2]{\6#1\9$^{#2}$}
\renewcommand{\ni}{\~{\i}}

\newcommand{\citepos}[1]{\citeauthor{#1}'s \citeyear{#1}}
\newcommand{\citeposs}[1]{\citeauthor{#1}'s}
\newcommand{\citetpos}[1]{\citeauthor{#1}'s (\citeyear{#1})}

\newcolumntype{R}[2]{%
    >{\adjustbox{angle=#1,lap=\width-(#2)}\bgroup}%
    l%
    <{\egroup}%
}
\newcommand*\rot{\multicolumn{1}{R{90}{0em}}}% no optional argument here, please!

\newcommand*\rots{\multicolumn{1}{R{90}{.7em}}}% no optional argument here, please!

% positive coefficients/difference
\definecolor{purple1}{RGB}{178,24,43}
\definecolor{purple2}{RGB}{239,138,98} 
\definecolor{purple3}{RGB}{253,219,199} 
%\definecolor{yellow4}{RGB}{255,255,204}
%\definecolor{green1}{RGB}{0, 158, 115} % >.95
%\definecolor{green2}{RGB}{55, 185, 141} % .85-.95
%\definecolor{green3}{RGB}{88, 214, 167} % .75-.85
%\definecolor{green4}{RGB}{119, 242, 194} % <.75

% negative coefficients/difference
\definecolor{yellow1}{RGB}{33,102,172}
\definecolor{yellow2}{RGB}{103,169,207}
\definecolor{yellow3}{RGB}{209,229,240}
%\definecolor{purple4}{RGB}{236,206,223}

\title{I don't know if naturalness ratings distinguish \\ presuppositions from nonpresuppositions. \\ Did Mandelkern et al.\ 2020 discover that they do?}

\author[$\circ$]{Judith Tonhauser}
\author[$\bullet$]{Judith Degen}
\affil[$\circ$]{University of Stuttgart, Department of Linguistics, Stuttgart, Germany, judith.tonhauser@ling.uni-stuttgart.de (corresponding author) \\ ORCID 0000-0002-6061-7120}
\affil[$\bullet$]{Prescient AI, New York, NY, United States, judithdegen@gmail.com \\ ORCID 0000-0003-2513-0234}

\renewcommand\Authands{ and }

\newcommand{\jt}[1]{\textbf{\color{blue}JT: #1}}

\begin{document}

\maketitle
\thispagestyle{empty}

\begin{abstract}

Presuppositions have long been taken to be a well-defined class of content that is characterized by typically projecting and by being conventionally associated with particular lexical items or constructions (e.g., \citealt{heim83,vds92,mandelkern-etal2020,presupposition-sep}). While the strength of the projection inference is often taken to distinguish presuppositions from nonpresuppositions, experimental investigations that used inference tasks suggested that the strength of the projection inference does not distinguish the two (e.g., \citealt{tbd-variability,white-rawlins-nels2018,demarneffe-etal-sub23,degen-tonhauser-language}). \citealt{mandelkern-etal2020} proposed that a different measure, namely comparison of naturalness ratings in explicit ignorance and support contexts, distinguishes presuppositions from nonpresuppositions. This paper presents the results of an experiment designed to investigate \citepos{mandelkern-etal2020} proposal in the domain of factive and nonfactive clause-embedding predicates. The results do not support \citepos{mandelkern-etal2020} proposal, which means that the question of whether presuppositional projection inferences can be distinguished from nonpresuppositional ones is still open.


\end{abstract}

\bigskip

\noindent
{\bf Keywords:} Projection inferences, presuppositions, inference ratings, naturalness ratings in explicit ignorance and support contexts.

\bigskip

\noindent
{\bf Acknowledgments:} 

\noindent
For helpful comments on this project we thank Craige Roberts, Gregory Scontras, Mandy Simons, the reviewers at \emph{Linguistics \& Philosophy}, and the audience at the syntax/semantics discussion group at the University of Stuttgart.

\bigskip

\noindent
{\bf Funding:} The experiment was funded by Judith Tonhauser's research funds.

\bigskip

\noindent
{\bf Competing interests:} The authors have no competing interests to declare.

\bigskip

\noindent
{\bf Ethics approval and informed consent:} The experiment was approved by the ethics review committee of the University of Stuttgart. Informed consent was obtained from the participants.

\bigskip

\noindent
{\bf Availability of materials, data, and code:} See footnote 7.

\bigskip

\noindent
{\bf Author contributions:} Both authors contributed to the conception and design of the experiment. Material preparation, data collection and analysis were performed by both authors. The first draft of the first and second versions of the manuscript were written by Judith Tonhauser. Both authors commented on previous versions of the manuscript. Both authors read and approved the final manuscript.


%\bibliographystyle{cslipubs-natbib}
\bibliographystyle{apacite}

%\bibliographystyle{/Users/tonhauser.1/Library/Latex/cslipubs-natbib}
\bibliography{bibliography}


 
\end{document}

